% 18-unsafe.tex — Clause 18: Unsafe code
\chapter{Unsafe code}
\label{sec:18-unsafe}
\index{unsafe code}

\section{Purpose}
\label{sec:unsafe.purpose}

\pnum
The \lain{unsafe} block provides a controlled escape hatch for operations
that a conforming implementation cannot statically verify.  Unsafe blocks
are explicitly delimited in the source text and are therefore auditable.
Unsafe code cannot accidentally spread outside the enclosing block.

\section{Syntax}
\label{sec:unsafe.syntax}

\pnum
An unsafe block is written as \lain{unsafe \{ body \}}.  It may appear
wherever a statement is valid (\S\,\ref{sec:stmt.unsafe}).

\section{What unsafe enables}
\label{sec:unsafe.enables}
\index{unsafe!what it enables}

\pnum
Inside an \lain{unsafe} block, the following operations are permitted:

\begin{enumerate}
  \item \textbf{Raw pointer dereference}: \lain{*ptr} reads the value
        pointed to by a raw pointer; \lain{*ptr = v} writes through a
        mutable raw pointer.
  \item \textbf{Address-of}: \lain{\&x} creates a raw pointer to the
        storage of \lain{x}.
  \item \textbf{Pointer and integer casts}: \lain{p as *T},
        \lain{n as *T}, \lain{p as usize}, etc.
  \item \textbf{Direct ADT field access}: accessing a field of an ADT
        variant without going through a \lain{case} arm.
  \item \textbf{Bounds-check suppression}: array indexing inside an
        \lain{unsafe} block does not require a VRA proof (the
        programmer asserts the access is in bounds).
\end{enumerate}

\section{What unsafe does NOT disable}
\label{sec:unsafe.not-disabled}
\index{unsafe!what it does not disable}

\pnum
The following safety features remain active inside \lain{unsafe} blocks:

\begin{longtable}{@{}ll@{}}
\toprule
\textbf{Feature} & \textbf{Disabled by \lain{unsafe}?} \\
\midrule
\endfirsthead
\toprule
\textbf{Feature} & \textbf{Disabled by \lain{unsafe}?} \\
\midrule
\endhead
\bottomrule
\endlastfoot
Pointer dereference                  & \textbf{Yes} \\
Pointer/integer casts                & \textbf{Yes} \\
Direct ADT field access              & \textbf{Yes} \\
Bounds check requirement             & \textbf{Yes} \\
Ownership tracking (\diagcode{E001}--\diagcode{E007}) & No \\
Borrow checking (\diagcode{E004})    & No \\
Move semantics                       & No \\
Type checking                        & No \\
Division-by-zero check (\diagcode{E015}) & No \\
Exhaustiveness checking (\diagcode{E014}) & No \\
Purity enforcement (\diagcode{E011}) & No \\
\end{longtable}

\begin{constraint}
  Ownership and linearity rules (\S\,\ref{sec:11-ownership}) are not
  suspended inside \lain{unsafe} blocks.  Use-after-move, double-move,
  and borrow conflicts are still constraint violations inside
  \lain{unsafe}~(\diagcode{E001}--\diagcode{E007}).
\end{constraint}

\section{Undefined behavior in unsafe}
\label{sec:unsafe.ub}
\index{undefined behavior!in unsafe}

\pnum
The following operations inside an \lain{unsafe} block may produce
\undefbehav{undefined behavior if the programmer's safety assertion is
incorrect}:

\begin{itemize}
  \item Dereferencing a null pointer.
  \item Dereferencing a pointer to freed or uninitialized memory.
  \item Writing to memory through a pointer derived from a read-only
        (shared) source.
  \item Accessing an array element with an out-of-range index.
  \item Accessing a field of an ADT through a tag that does not match
        the active variant.
\end{itemize}

\pnum
A conforming implementation that generates C99 translation output inherits
the undefined behavior rules of ISO/IEC~9899 (C17) for operations within
\lain{unsafe} blocks.

\section{Borrow scope for unsafe blocks}
\label{sec:unsafe.borrow-scope}

\pnum
An \lain{unsafe} block is a regular block scope for the purposes of borrow
checking.  Borrows registered within the block are released on block exit,
and the borrow checker's invariants hold at the block boundaries.

\section{Nesting}
\label{sec:unsafe.nesting}

\pnum
\lain{unsafe} blocks may be nested.  The inner block does not extend the
permissions of the outer block beyond what is already enabled.
\idbehav{Whether the implementation issues a diagnostic for redundant nested
unsafe blocks is implementation-defined}.
